\section*{Цель работы}
Исследование и определение ключевых метрологических характеристик электронных вольтметров: основной погрешности, амплитудно-частотной характеристики (АЧХ) и влияния формы измеряемого напряжения.

\section*{Задачи работы}
\begin{enumerate}
    \item Ознакомиться с экспериментальной установкой и методикой проведения измерений.
    \item Определить основную погрешность поверяемого электронного вольтметра методом сличения с образцовым прибором. Построить график зависимости относительной и приведенной погрешностей от показаний вольтметра и сделать вывод о его соответствии заявленному классу точности.
    \item Определить амплитудно-частотную характеристику (АЧХ) вольтметра. Построить график АЧХ и найти рабочую полосу частот прибора.
    \item Экспериментально исследовать влияние формы измеряемого напряжения (синусоидальной, прямоугольной, треугольной) на показания вольтметра.
\end{enumerate}

\section*{Теоретические сведения}

\subsection*{Определение основной погрешности}
Основную погрешность электронного вольтметра определяют в нормальных условиях методом сличения, сравнивая его показания $U$ с показаниями образцового (более точного) прибора $U_0$. В данной работе в качестве образцового прибора используется цифровой вольтметр.

Ключевыми показателями погрешности являются:
\begin{itemize}
    \item \textit{Абсолютная погрешность} — разница между показанием прибора и действительным значением измеряемой величины.
    \begin{equation}
        \label{eq:abs_err}
        \Delta U = U - U_0
    \end{equation}
    \item \textit{Относительная погрешность} — отношение абсолютной погрешности к действительному значению.
    \begin{equation}
        \label{eq:rel_err}
        \delta = \frac{\Delta U}{U_0} \cdot 100\%
    \end{equation}
    \item \textit{Приведенная погрешность} — отношение абсолютной погрешности к нормирующему значению $U_N$ (обычно верхний предел измерений шкалы).
    \begin{equation}
        \label{eq:reduced_err}
        \gamma = \frac{\Delta U}{U_N} \cdot 100\%
    \end{equation}
    \item \textit{Вариация показаний} — разница между показаниями прибора в одной и той же точке, полученными при подходе к ней со стороны меньших и больших значений.
    \begin{equation}
        \label{eq:variation}
        H = |U_{\text{о.ув}} - U_{\text{о.ум}}|
    \end{equation}
    где $U_{\text{о.ув}}$ и $U_{\text{о.ум}}$ — показания образцового вольтметра при увеличении и уменьшении сигнала соответственно.
\end{itemize}

\subsection*{Определение амплитудно-частотной характеристики (АЧХ)}
АЧХ описывает зависимость показаний вольтметра от частоты входного синусоидального сигнала при неизменной амплитуде. Идеальный вольтметр имеет равномерную АЧХ в рабочем диапазоне частот.

\textit{Рабочая полоса частот} — это диапазон $[f_\text{н}, f_\text{в}]$, в пределах которого неравномерность АЧХ не превышает установленного допуска. Для вольтметра GVT-417B допускается спад АЧХ не более чем на $10\%$ относительно показания на опорной частоте $f_0 = 1~\text{кГц}$.

Коэффициент передачи (нормированная АЧХ) вычисляется по формуле:
\begin{equation}
    \label{eq:afc_coeff}
    K(f) = \frac{U(f)}{U(f_0)}
\end{equation}
где $U(f)$ — показание вольтметра на частоте $f$, а $U(f_0)$ — показание на опорной частоте $f_0$. Граничные частоты $f_\text{н}$ и $f_\text{в}$ определяются из условия $K(f) = 0.9$.

\subsection*{Влияние формы входного сигнала}
Электронные вольтметры переменного тока градуируются в действующих (среднеквадратичных) значениях синусоидального напряжения. При измерении несинусоидальных сигналов возникают дополнительные погрешности.

Вольтметр GVT-417B оснащен преобразователем средневыпрямленного значения. Его показания $U_\text{п}$ пропорциональны среднему значению выпрямленного сигнала $U_\text{ср}$. Связь между действующим значением $U$ и средневыпрямленным $U_\text{ср}$ определяется \textit{коэффициентом формы} $\text{к}_\text{ф} = U/U_\text{ср}$.
\begin{equation}
    \label{eq:form_factor_relation}
    U = \text{к}_\text{ф} U_\text{ср}
\end{equation}
Для синусоидального сигнала $\text{к}_\text{ф} \approx 1.11$. Шкала прибора откалибрована с учетом этого коэффициента. Для сигналов другой формы:
\begin{itemize}
    \item Прямоугольный сигнал (меандр): $\text{к}_\text{ф} = 1.0$.
    \item Треугольный сигнал: $\text{к}_\text{ф} \approx 1.15$.
\end{itemize}
Дополнительная относительная погрешность $\delta_\text{ф}$ из-за несоответствия формы сигнала определяется как:
\begin{equation}
    \label{eq:waveform_err}
    \delta_\text{ф} = \frac{U_\text{п} - U}{U} \cdot 100\%
\end{equation}
где $U_\text{п}$ — показания вольтметра, а $U$ — истинное действующее значение сигнала.

\section*{Порядок выполнения работы}

\subsection*{Схема установки}
Для проведения измерений используется схема, приведенная на \cref{fig:setup_scheme}.
\begin{figure}[H]
    \centering
    \includegraphics[width=0.8\linewidth]{figs/setup_scheme.pdf}
    \caption{Схема экспериментальной установки: ГС — генератор сигналов, ЦВ — цифровой вольтметр (образцовый), ЭВ — электронный вольтметр (поверяемый), ЭЛО — электронно-лучевой осциллограф}
    \label{fig:setup_scheme}
\end{figure}

\subsection*{Задание 2. Определение основной погрешности}
\begin{enumerate}
    \item Соберите установку в соответствии со схемой на \cref{fig:setup_scheme}.
    \item На генераторе сигналов (ГС) установите синусоидальный сигнал с частотой $f_0 = 1~\text{кГц}$.
    \item Выберите на поверяемом вольтметре (ЭВ) предел измерений, указанный преподавателем (например, $1~\text{В}$ или $3~\text{В}$). Запишите нормирующее значение $U_N$ в \cref{tab:error_measurement}.
    \item Выберите $6-10$ равномерно распределенных по шкале отметок $U_i$ для поверки.
    \item Для каждой отметки $U_i$:
    \begin{itemize}
        \item Плавно увеличивая напряжение на выходе ГС, добейтесь, чтобы стрелка ЭВ точно совпала с отметкой $U_i$. Зафиксируйте показание образцового вольтметра (ЦВ) $U_{\text{о.ув},i}$ и занесите его в \cref{tab:error_measurement}.
        \item Плавно уменьшая напряжение (от значения, превышающего $U_i$), снова добейтесь совпадения стрелки ЭВ с отметкой $U_i$. Зафиксируйте показание ЦВ $U_{\text{о.ум},i}$ и занесите его в \cref{tab:error_measurement}.
    \end{itemize}
\end{enumerate}

\subsection*{Задание 3. Определение АЧХ}
\begin{enumerate}
    \item На ГС установите синусоидальный сигнал с частотой $f_0 = 1~\text{кГц}$.
    \item С помощью регулятора выходного напряжения ГС установите на ЭВ показание $U(f_0)$, равное $(0.7 \dots 0.9) \cdot U_N$. Запишите это значение в \cref{tab:afc_high_freq} и \cref{tab:afc_low_freq}.
    \item Для контроля постоянства амплитуды выходного сигнала используйте ЭЛО. Зафиксируйте амплитуду (или размах) осциллограммы. \textit{Важно: в ходе всего эксперимента по снятию АЧХ амплитуда сигнала на выходе ГС должна оставаться неизменной.}
    \item \textit{Определение верхней граничной частоты $f_\text{в}$}:
    \begin{itemize}
        \item Увеличивайте частоту на ГС с шагом, указанным в \cref{tab:afc_high_freq}. Для каждой частоты $f_i$ записывайте показание ЭВ $U(f_i)$.
        \item В области, где показания начинают заметно падать (приближаются к $0.9 \cdot U(f_0)$), уменьшите шаг изменения частоты для более точного определения $f_\text{в}$.
    \end{itemize}
    \item \textit{Определение нижней граничной частоты $f_\text{н}$}:
    \begin{itemize}
        \item Вернитесь к частоте $f_0 = 1~\text{кГц}$ и проверьте исходное значение $U(f_0)$.
        \item Уменьшайте частоту на ГС с шагом, указанным в \cref{tab:afc_low_freq}. Для каждой частоты $f_i$ записывайте показание ЭВ $U(f_i)$.
        \item В области спада показаний уменьшите шаг для точного определения $f_\text{н}$.
    \end{itemize}
\end{enumerate}

\subsection*{Задание 4. Определение влияния формы сигнала}
\begin{enumerate}
    \item На ГС установите синусоидальный сигнал с частотой, лежащей в рабочей полосе частот вольтметра (например, $f = 1~\text{кГц}$).
    \item Установите напряжение на выходе ГС так, чтобы показание ЭВ $U_\text{п}$ составляло $(0.5 \dots 0.6) \cdot U_N$. Запишите это значение в \cref{tab:waveform_influence} для синусоидальной формы.
    \item \textit{Важно: не изменяйте положение регулятора амплитуды на ГС до конца эксперимента.}
    \item Переключите форму сигнала на ГС на прямоугольную. Запишите новое показание ЭВ $U_\text{п}$ в \cref{tab:waveform_influence}.
    \item Переключите форму сигнала на ГС на треугольную. Запишите показание ЭВ $U_\text{п}$ в \cref{tab:waveform_influence}.
\end{enumerate}

\section*{Обработка результатов}
\begin{enumerate}
    \item Для \textit{задания 2} рассчитать абсолютную, относительную, приведенную погрешности и вариацию для каждого измерения, используя формулы \cref{eq:abs_err,eq:rel_err,eq:reduced_err,eq:variation}. Построить на одном графике зависимости $\delta(U_0)$ и $\gamma(U_0)$.
    \item Для \textit{задания 3} рассчитать коэффициент $K(f)$ по формуле \cref{eq:afc_coeff} для всех частот. Построить единый график АЧХ $K(f)$ в логарифмическом масштабе по оси частот. Определить по графику граничные частоты $f_\text{н}$ и $f_\text{в}$ на уровне $0.9$.
    \item Для \textit{задания 4} рассчитать теоретические показания вольтметра для прямоугольного и треугольного сигналов и вычислить дополнительную погрешность по \cref{eq:waveform_err}. Сравнить расчетные и экспериментальные данные.
\end{enumerate}

% Изображение \cref{fig:setup_scheme} соответствует рис. 2.1 из методического пособия.